\documentclass[11pt,a4paper]{article}
\usepackage[margin=27mm]{geometry}
\usepackage{amsmath,amssymb,amsthm,booktabs,hyperref}
\hypersetup{hidelinks}
\setlength{\emergencystretch}{1em}
\newtheorem{theorem}{Theorem}[section]
\newtheorem{lemma}[theorem]{Lemma}
\newtheorem{proposition}[theorem]{Proposition}
\newtheorem{corollary}[theorem]{Corollary}
\numberwithin{equation}{section}
\newcommand{\im}{\operatorname{im}}
\newcommand{\rank}{\operatorname{rank}}
\newcommand{\wt}{\operatorname{wt}}
\title{Incidence operators, orthogonal codes and permutation modules}
\author{Common theory for two applications}
\date{}
\begin{document}\maketitle

The starting point is a symmetric operator on a space with its coordinate
bilinear form. An orthogonal decomposition of its image can often be
read from a short polynomial identity, or transferred from a much smaller
incidence space. This separates the questions of dimension and duality
from the questions of minimum weight and group action. The latter require
additional arguments: averaging over coordinate permutations reduces a
distance search, a matrix algebra certifies irreducibility, and refinement
of a support incidence structure bounds the full automorphism group.
Each argument is proved once here and then applied to literal matrices
and finite data in the accompanying papers.

Our two source manuscripts are Leemans and Rodrigues,
\emph{Linear codes with complementary duals from some strongly regular
subgraphs of the McLaughlin graph}, Math. Commun.\ 21 (2016), 239--249
(L050), and Key and Rodrigues, \emph{LCD codes from adjacency matrices of
graphs}, the seventeen-page manuscript dated 4 August 2017 (F017).
The graph constructions, original code assertions, and results quoted
from their predecessors remain attributed to those sources in the
applications. Orthogonal projections, incidence counting, the simplex
chain contraction, MacWilliams duality and elementary module theory are
classical tools. Our contribution here is their common derivation and
the explicitly checked passage to the stated applications. Historical
priority for these methods is not asserted.

The common complement theorem includes the possible one-dimensional
defect and does not require an integral spectrum. The subset and Hamming
calculations give families of consequences of the same operator tests.
The binary incidence argument also determines all minimum words, whereas
the ternary triple-incidence calculation corrects a dimension in F017.
Other corrections, retained classification results and finite computation
premises are stated with the exact source assertions in each application.
No numerical compression factor is assigned by this draft.

The whole of this common paper is a prerequisite for the applications
when lengths are compared. A theorem used by only one application is
still charged here in full. Executable certificates have a separate
inventory; any manuscript comparison excludes executable programs on
both sides. The nine generic Lean declarations already associated with
the projection and averaging arguments retain their limited stated
scope. They do not formalize all the proofs in this paper.

\section{Symmetric operators and graph duality}

Let $F$ be a field and let $V=F^n$, where $n>0$. We use the standard
nondegenerate symmetric bilinear form $\langle x,y\rangle=x^{\mathsf T}y$.
For a symmetric matrix $A$, write $C_A=\im A$. Thus $C_A^\perp=\ker A$.
The code $C_A$ is LCD if $C_A\cap C_A^\perp=0$.
Write $j=(1,\ldots,1)^{\mathsf T}$ and $J=jj^{\mathsf T}$; integers in matrix
identities denote their images in $F$.

\begin{lemma}\label{ULCD.Projector}
If $E$ is symmetric and $E^2=E$, then
$V=\im E\mathbin{\perp\!\oplus}\ker E$.
In particular $\im E$ is LCD. If $E_1,\ldots,E_t$ are symmetric idempotents
with $E_iE_k=0$ for $i\ne k$ and $\sum E_i=I$, their images form an
orthogonal direct sum equal to $V$.
\end{lemma}
\begin{proof}
Every $v$ is $Ev+(v-Ev)$, with the two terms in the asserted spaces.
If $x=Ey$ and $Ex=0$, then $x=E^2y=Ex=0$. Moreover
$\langle Ey,z\rangle=\langle y,Ez\rangle=0$ when $Ez=0$.
For several idempotents, $v=\sum E_iv$ gives spanning, and application of
$E_k$ to a relation $\sum x_i=0$, $x_i\in\im E_i$, gives $x_k=0$.
The same inner-product identity proves pairwise orthogonality.
\end{proof}

\begin{theorem}\label{ULCD.Complement}
Let $A$ be symmetric and put $B=J-I-A$. Suppose that $AB=0$.
There is a scalar $\nu\in F$ for which $Aj=\nu j$.
Put $\bar\nu=n-1-\nu$. Then $\nu\bar\nu=0$, both $C_A$ and $C_B$
are LCD, and the following are orthogonal decompositions with the indicated
projection matrices:
\[
\begin{array}{c|c|c}
\text{condition}&\text{decomposition}&\text{projections}\\\hline
\nu\ne0&V=C_A\perp\!\oplus C_B&J-A,\ -B\\
\bar\nu\ne0&V=C_A\perp\!\oplus C_B&-A,\ J-B\\
\nu=\bar\nu=0&V=C_A\perp\!\oplus C_B\perp\!\oplus Fj&-A,\ -B,\ J.
\end{array}
\]
Consequently $C_A=C_B^\perp$ if and only if $(\nu,\bar\nu)\ne(0,0)$.
In the exceptional case $C_B^\perp=C_A\perp\!\oplus Fj$ and
$C_A^\perp=C_B\perp\!\oplus Fj$.
\end{theorem}
\begin{proof}
Transposition gives $BA=0$. Expanding these two products gives
\[
 AJ=A^2+A=JA.
\]
If $d=Aj$, then $AJ=dj^{\mathsf T}$ and $JA=jd^{\mathsf T}$.
Equality of their $(r,s)$ entries says $d_r=d_s$ for all $r,s$.
Since $n>0$, this defines $\nu$, and hence $Bj=\bar\nu j$.
Applying $AB=0$ to $j$ gives $\nu\bar\nu=0$. We have also proved
\[
 A^2+A=\nu J,\quad B^2+B=\bar\nu J,\quad
 AJ=JA=\nu J,\quad BJ=JB=\bar\nu J,\quad J^2=nJ.
\]
Suppose first that $\nu\ne0$. Then $\bar\nu=0$ and $n-\nu=1$.
Set $E=J-A$ and $F=-B$. Direct multiplication gives
\[
 E^2=(n-\nu)J-A=E,\quad F^2=B^2=-B=F,\quad
 EF=FE=0,\quad E+F=I.
\]
Here $\im F=\im B$. Since $Aj=\nu j$ with $\nu\ne0$, we have
$j\in\im A$, so $\im(J-A)\subseteq\im A$.
Conversely $EA=JA-A^2=A$, so $\im A\subseteq\im E$.
Lemma~\ref{ULCD.Projector} proves the first row. Interchanging $A$ and $B$
proves the second.

Finally, if $\nu=\bar\nu=0$, then $n=1$ in $F$.
The displayed identities show that $-A,-B,J$ are symmetric, pairwise
annihilating idempotents, with sum $I$. Their images are $C_A,C_B,Fj$.
Since $j\ne0$ and $\langle j,j\rangle=n=1$, the last summand is a
nonzero nondegenerate line. This proves the third row and both descriptions
of the orthogonal complements. In particular that line prevents equality
$C_A=C_B^\perp$ in precisely this case.
\end{proof}

\begin{corollary}\label{ULCD.Graph}
Let $A$ be the adjacency matrix of a nonempty finite simple graph over $F$
and $B$ that of its complement. If the two graph codes are orthogonal,
then both are LCD. The vertices have equal degrees in $F$, even when their
integer degrees are not equal. The complement code is the full orthogonal
dual unless both degrees vanish in $F$, in which case the missing summand
is exactly the constant line.
\end{corollary}
\begin{proof}
Orthogonality of the column spaces is $A^{\mathsf T}B=AB=0$.
Apply Theorem~\ref{ULCD.Complement}. Matrix multiplication by $j$ records
the vertex degrees.
\end{proof}

\begin{proposition}\label{ULCD.Reflexive}
If $A$ is symmetric and $A(A+I)=0$, then $-A$ and $A+I$ are the
orthogonal projections onto $C_A$ and $C_{A+I}$, respectively. Thus
$C_A=C_{A+I}^{\perp}$. Conversely this equality implies $A(A+I)=0$.
\end{proposition}
\begin{proof}
The two matrices are symmetric, sum to $I$, annihilate one another, and
are idempotent because $A^2=-A$. Their images are as stated, so
Lemma~\ref{ULCD.Projector} applies. For the converse, the asserted equality
implies orthogonality of all columns, which gives $A(A+I)=0$.
\end{proof}

\begin{proposition}\label{ULCD.SRG}
Let $A$ be the adjacency matrix of a noncomplete, nonempty strongly regular
graph with parameters $(v,k,\lambda,\mu)$, having both an edge and a
nonedge. Over a field of characteristic $p$, it satisfies
\[
 A(A+I)=0\quad\Longleftrightarrow\quad
 k=0,\quad\lambda=-1,\quad\mu=0 \pmod p,
\]
and
\[
 A(J-I-A)=0\quad\Longleftrightarrow\quad
 k-\mu=0,\quad k-1-\lambda=0\pmod p.
\]
The first condition gives the reflexive-code decomposition of
Proposition~\ref{ULCD.Reflexive}. The second gives the exact decomposition
of Theorem~\ref{ULCD.Complement}, with $\nu=k$ and $\bar\nu=v-1-k$.
\end{proposition}
\begin{proof}
Counting common neighbours gives
$A^2=(k-\mu)I+(\lambda-\mu)A+\mu J$.
Its diagonal, edge and nonedge entries give the first equivalence directly.
For the second, the edge entries of $A(J-I-A)$ equal $k-1-\lambda$,
the nonedge entries equal $k-\mu$, and the diagonal entries are zero.
All three types occur under the stated assumptions.
\end{proof}

\section{Incidence transfer and residual designs}

\begin{lemma}\label{ULCD.IncidenceTransfer}
Let $H$ be an $m\times n$ matrix over a field, with nonsingular Gram matrix
$G=HH^{\mathsf T}$. Let $A$ be symmetric and suppose
\[
 A=H^{\mathsf T}MH
\]
for an $m\times m$ matrix $M$. Then
\[
 \im A=H^{\mathsf T}(\im M),\qquad \rank A=\rank M.
\]
Under this identification the form on $\im A$ is the restriction to
$\im M$ of $(x,y)\mapsto x^{\mathsf T}Gy$.
Consequently $\im A$ is LCD exactly when this restricted form is
nondegenerate.
\end{lemma}
\begin{proof}
Nonsingularity of $G$ implies that $H$ is onto and $H^{\mathsf T}$ is
injective: $H(H^{\mathsf T}G^{-1}x)=x$, and $H^{\mathsf T}x=0$ implies
$Gx=0$. Hence $MH$ has image $\im M$, and applying the injective map
$H^{\mathsf T}$ proves the first two assertions. For $x,y\in\im M$,
\[
 \langle H^{\mathsf T}x,H^{\mathsf T}y\rangle=x^{\mathsf T}Gy.
\]
The radical of the restricted form on a code is its intersection with its
orthogonal complement. This proves the last assertion.
\end{proof}

Let $\mathcal D$ be a Steiner system $S(4,7,23)$. Fix a set $T$ of size
$t\in\{0,1,2\}$. Retain the heptads disjoint from $T$, and write
$H_t$ for their point--block incidence matrix on the remaining $23-t$
points. Two distinct heptads of $\mathcal D$ meet in either one or three
points. This intersection property is a retained property of this Steiner
system; the explicit binary construction used in our finite certificate
also verifies it for every pair. Define $A_t$ on the retained blocks by
joining two distinct blocks precisely when they meet in three points.
All matrix calculations in the next theorem take place over $\mathbb F_3$.
Write $e$ for the constant vector on points and $j$ for the constant vector
on blocks; the subscripts on identity and all-one matrices indicate their
sizes.

\begin{theorem}\label{ULCD.ResidualSteiner}
The three graph codes have the following descriptions, dimensions and
orthogonal decompositions:
\[
\begin{array}{c|c|c|c|c}
t&n& H_tH_t^{\mathsf T}&\im A_t&\rank A_t\\\hline
0&253&2I_{23}&\im H_0^{\mathsf T}&23\\
1&176&I_{22}+J_{22}&H_1^{\mathsf T}(e^\perp)&21\\
2&120&I_{21}&\im H_2^{\mathsf T}&21.
\end{array}
\]
Each of these codes is LCD. Moreover, if $B_t=J_n-I_n-A_t$ is the
complementary adjacency matrix, then
\[
 (\im A_t)^\perp=\im B_t\quad(t=1,2).
\]
For $t=0$, the code has the orthogonal splitting
\[
 \im A_0=\mathbb F_3j\perp\!\oplus H_0^{\mathsf T}(e^\perp),
\]
whose second summand has dimension $22$ and is LCD. The orthogonal dual
of $\im A_t$ has dimension $n-\rank A_t$ in each case.
\end{theorem}
\begin{proof}
The numbers of heptads through an $s$-set, for $0\leq s\leq4$, are
\[
 \lambda_s=\frac{\binom{23-s}{4-s}}{\binom{7-s}{4-s}}
       \quad\text{and hence}\quad
 (\lambda_0,\lambda_1,\lambda_2,\lambda_3,\lambda_4)
       =(253,77,21,5,1).
\]
Indeed, count the four-subsets containing that $s$-set and use the unique
heptad containing each four-subset. Inclusion--exclusion on the points of
$T$ now gives the number $n_t$ of retained blocks, the number $r_t$ through
each remaining point and the number $s_t$ through each remaining pair:
\[
\begin{array}{c|rrr}
t&n_t&r_t&s_t\\\hline
0&253&77&21\\
1&176&56&16\\
2&120&40&12.
\end{array}
\]
Thus $H_tH_t^{\mathsf T}=(r_t-s_t)I+s_tJ$, giving the three Gram matrices
in the statement. They are nonsingular. Only the middle one needs comment:
if $(I+J)x=0$ and $u=e^{\mathsf T}x$, then $x=-ue$ and
$u=-22u=-u$ in $\mathbb F_3$, whence $u=0$ and $x=0$.

Every column of $H_t$ has sum $7=1$ in $\mathbb F_3$, so
$H_t^{\mathsf T}e=j$. The diagonal entries of $H_t^{\mathsf T}H_t$ are
also $1$. Off the diagonal they are $1$ at intersection size one and $0$
at intersection size three. Consequently
\[
 A_t=J_n-H_t^{\mathsf T}H_t
     =H_t^{\mathsf T}(J_{23-t}-I_{23-t})H_t.                 \tag{2.1}\label{ULCD.ResidualIdentity}
\]
For $m=23$ or $21$, the matrix $J_m-I_m$ is invertible over $\mathbb F_3$.
In fact, $(J_m-I_m)x=0$ implies $x=ue$ and $u=mu$; since $m-1\ne0$
in both cases, $u=0$. For $m=22$, its kernel is $\mathbb F_3e$, its image
is $e^\perp$, and it acts as $-I$ on that image. The image assertion follows
from $e^{\mathsf T}(J_m-I_m)=0$ and the rank $m-1$.
Lemma~\ref{ULCD.IncidenceTransfer} proves every image and dimension in the table.

For $t=0,2$, the image in the point space is the whole space and its Gram
matrix is nonsingular, proving LCD. For $t=1$, $G=I+J$ restricts to the
ordinary dot product on $e^\perp$. This restriction is nondegenerate
because $e^{\mathsf T}e=22=1$: a vector in its radical belongs both to
$e^\perp$ and to $\mathbb F_3e$, so vanishes. For $t=0$ the same argument
uses $G=2I$ and $e^{\mathsf T}e=23=2$ to give the asserted splitting.
Its constant line is nondegenerate since $j^{\mathsf T}j=253=1$.

It remains to identify the complementary code when $t=1,2$.
By \eqref{ULCD.ResidualIdentity}, $B_t=H_t^{\mathsf T}H_t-I_n$. Put $M=J_{23-t}-I_{23-t}$ and
$G=H_tH_t^{\mathsf T}$. Then
\[
 A_tB_t=H_t^{\mathsf T}M(G-I)H_t=0:
\]
for $t=2$, $G=I$; for $t=1$, $G-I=J_{22}$ and
$(J_{22}-I_{22})J_{22}=(22-1)J_{22}=0$. In the latter case
$B_tH_t^{\mathsf T}e=H_t^{\mathsf T}(G-I)e=H_t^{\mathsf T}e=j$;
in the former case $A_tj=2j$, since \eqref{ULCD.ResidualIdentity} and $H_tj=r_te=e$ give
$A_tj=H_t^{\mathsf T}(21-1)e=2j$.
Thus in both cases at least one graph degree is nonzero in the field.
Theorem~\ref{ULCD.Complement} yields
$\im B_t=(\im A_t)^\perp$. 
\end{proof}

\begin{proposition}\label{ULCD.SteinerRecognition}
Suppose, in addition, that the length-$253$ code in
Theorem~\ref{ULCD.ResidualSteiner} has minimum weight $77$ and exactly $46$ words of
that weight. Then its group of coordinate permutations preserving the code
is exactly the group induced on the blocks by $\operatorname{Aut}(\mathcal D)$.
\end{proposition}
\begin{proof}
For each of the $23$ points $p$, its incidence row $h_p$ is a codeword of
weight $77$. The $46$ words $h_p,-h_p$ are distinct: the incidence rows
are independent by the nonsingular Gram matrix, and a nonzero row with
entries in $\{0,1\}$ cannot equal the negative of another such row over
$\mathbb F_3$. The hypothesis therefore lists all minimum words.
A coordinate permutation preserving the code must permute their $23$
supports. These supports are exactly the sets of blocks through the
respective points. The induced permutation of points and the given
permutation of blocks preserve incidence, so form an automorphism of
$\mathcal D$. Conversely every design automorphism permutes the incidence
rows and hence preserves their span, which is the code by
Theorem~\ref{ULCD.ResidualSteiner}. This proves equality of the induced groups.
\end{proof}

The minimum-word hypothesis in Proposition~\ref{ULCD.SteinerRecognition}
requires its own exact certificate. The application also identifies the
full design group by literal permutations attaining an incidence-refinement
bound, using Theorem~\ref{ULCD.Refinement}. Neither incidence transfer nor
recognition alone establishes irreducibility of the augmentation summand.

\section{Binary incidence spaces}

Let $\Omega$ be an $n$-element set with $n>0$. Let $V_k$ be the vector space over
$\mathbb F_2$ with orthonormal basis $e_S$ indexed by the $k$-subsets of
$\Omega$, including $V_0=\mathbb F_2e_\varnothing$. Put $V_{-1}=V_{n+1}=0$.
Define
\[
 d_k e_S=\sum_{x\in S}e_{S\setminus\{x\}},\qquad
 u_k e_S=\sum_{x\notin S}e_{S\cup\{x\}}.
\]
Thus $u_{k-1}=d_k^{\mathsf T}$. The space
$W_{k-1}=\im u_{k-1}\leq V_k$ is generated by the indicator vectors of
all $k$-sets containing a fixed $(k-1)$-set, exactly the binary incidence
code denoted by $W_{k-1}$ in F017. Its dual is $\ker d_k$.

\begin{lemma}\label{ULCD.Simplex}
The following identities hold, with all scalar integers reduced modulo two:
\[
 d_{k-1}d_k=0,\qquad u_{k+1}u_k=0,\qquad
 d_{k+1}u_k+u_{k-1}d_k=nI_{V_k}.
\]
The complex is exact. For $1\leq k\leq n$,
\[
 \rank d_k=\binom{n-1}{k-1},\qquad
 \dim\ker d_k=\binom{n-1}{k}.
\]
\end{lemma}
\begin{proof}
Each $(k-2)$-subset of a fixed $k$-set is obtained by two orders of deletion,
so $d^2=0$. Transposition gives $u^2=0$. In $d_{k+1}u_ke_S$, the coefficient
of $e_S$ is $n-k$, while in $u_{k-1}d_ke_S$ it is $k$. Each other possible
term replaces one element of $S$ by an element outside $S$; it occurs once
in each expression and cancels. This proves the third identity.

For exactness, fix $v\in\Omega$ and define $h_ke_S=e_{S\cup\{v\}}$ if
$v\notin S$, and $h_ke_S=0$ otherwise. Directly on a basis vector,
\[
 d_{k+1}h_k+h_{k-1}d_k=I_{V_k}.
\]
If $v\notin S$, deleting the newly added $v$ contributes $e_S$, and all
other terms occur twice. If $v\in S$, only deletion and subsequent addition
of $v$ survives. Thus if $d_kc=0$, then $c=d_{k+1}h_kc$, proving
$\ker d_k=\im d_{k+1}$. Now $\rank d_1=1$, since every singleton maps
to $e_\varnothing$. Exactness and rank--nullity give
$\rank d_{k+1}=\binom nk-\rank d_k$; induction with Pascal's identity
proves both dimension formulas. Transposing also gives
$\ker u_k=\im u_{k-1}=W_{k-1}$.
\end{proof}

\begin{theorem}\label{ULCD.Johnson}
Let $1\leq k<n$, let $A$ be the binary adjacency matrix of $J(n,k)$, and
put $E=u_{k-1}d_k$. Then
\[
 E^2=nE,\qquad A=E+(k\bmod2)I.
\]
If $n$ is odd, $E$ is the orthogonal projector onto $W_{k-1}$. Consequently
\[
\begin{array}{c|cc}
 &\im A&\im(A+I)\\\hline
k\text{ even}&W_{k-1}&W_{k-1}^{\perp}\\
k\text{ odd}&W_{k-1}^{\perp}&W_{k-1}.
\end{array}
\]
Both codes are LCD. If $n$ is even, $E^2=0$ and
$\im E\subseteq W_{k-1}\cap W_{k-1}^{\perp}$. In that case
\[
\begin{array}{c|cc}
 &\im A&\im(A+I)\\\hline
k\text{ even}&\im E&V_k\\
k\text{ odd}&V_k&\im E.
\end{array}
\]
\end{theorem}
\begin{proof}
The diagonal entry of $E=d_k^{\mathsf T}d_k$ is $k$. For two distinct
$k$-sets, the entry counts common $(k-1)$-subsets, so it is one precisely
when their intersection has size $k-1$. This proves the formula for $A$.
Using Lemma~\ref{ULCD.Simplex} in degree $k-1$ gives
\[
 E^2=u_{k-1}(nI-u_{k-2}d_{k-1})d_k=nE,
\]
and the same identities give
$E u_{k-1}=n u_{k-1}$ and $d_kE=n d_k$.
If $n$ is odd, the first identity shows that $E$ acts as the identity on
$W_{k-1}$, while its definition shows that its image is contained there.
Since $E$ is symmetric and idempotent, its kernel is the orthogonal
complement of its image, and $I+E$ projects onto that kernel.
The table for odd $n$ follows.
If $n$ is even, $\im E\subseteq\im u_{k-1}=W_{k-1}$ and $d_kE=0$, giving
the hull inclusion. Finally $(I+E)^2=I$, so $I+E$ is invertible, which proves
the second table.
\end{proof}

\begin{theorem}\label{ULCD.SimplexMin}
For $1\leq k<n$, the binary code $W_{k-1}$ has minimum distance $n-k+1$.
Its minimum words are exactly
\[
 u_{k-1}e_R=\sum_{S\supset R,\ |S|=k}e_S,
 \qquad |R|=k-1.
\]
The dual code $\ker d_k$ has minimum distance $k+1$. Its minimum words
are exactly
\[
 d_{k+1}e_T=\sum_{S\subset T,\ |S|=k}e_S,
 \qquad |T|=k+1.
\]
\end{theorem}
\begin{proof}
Let $0\ne c\in W_{k-1}=\ker u_k$, and choose $K$ in its support. For
each $x\notin K$, the coefficient of $e_{K\cup\{x\}}$ in $u_kc$ is zero.
Since $K$ contributes one, at least one other set $K\setminus\{i\}\cup\{x\}$
belongs to the support. These choices are distinct for different $x$, so
$\wt(c)\geq1+(n-k)$. The proposed star words attain this bound.

Suppose equality holds. For each $x\notin K$ there is exactly one chosen
$i=f(x)\in K$, and these $n-k$ sets together with $K$ comprise the entire
support. If there are at least two elements $x,y\notin K$, apply $u_kc=0$
to the $(k+1)$-set $K\setminus\{i\}\cup\{x,y\}$ for each $i\in K$.
The only possible supported facets are
$K\setminus\{i\}\cup\{x\}$ and $K\setminus\{i\}\cup\{y\}$: every other
facet contains both outside elements. Hence $f(x)=i$ if and only if
$f(y)=i$, and $f$ is constant. If there is only one outside element,
$f$ is already constant. The support is therefore the star of
$K\setminus\{f(x)\}$, as asserted. This also covers $k=1$.

Now let $0\ne c\in\ker d_k$ and choose $K$ in its support. For every
$i\in K$, the coefficient of $e_{K\setminus\{i\}}$ in $d_kc$ forces a
second supported set containing that facet. Such a set has the form
$K\setminus\{i\}\cup\{x_i\}$ with $x_i\notin K$. For distinct $i$ these
sets are distinct, so $\wt(c)\geq k+1$. The proposed boundary words attain
the bound, since $d_kd_{k+1}=0$.

At equality, those $k$ sets and $K$ are the full support. If $k\geq2$,
consider distinct $i,j\in K$ and the $(k-1)$-set
$K\setminus\{i,j\}\cup\{x_i\}$. It lies in the supported set
$K\setminus\{i\}\cup\{x_i\}$. The only possible second supported set
containing it is $K\setminus\{j\}\cup\{x_j\}$, and containment forces
$x_i=x_j$. Thus all $x_i$ are the same element $x$, and the support is the
boundary of $K\cup\{x\}$. If $k=1$, every word of weight two is already
the boundary of a two-element set. This proves exhaustion.
\end{proof}

\section{Ternary incidence on triples}

Let $\Omega$ have $n$ elements, and let $V$ be the vector space over
$\mathbb F_3$ with coordinates indexed by the triples of $\Omega$.
For each pair $e\subset\Omega$, write $v_e\in V$ for the indicator of
the triples containing $e$. Let $W_2$ be the span of these vectors and
let $\mathbf j$ be the constant vector. All orthogonal complements below
refer to the usual coordinate inner product.

\begin{lemma}\label{ULCD.TernaryIncidence}
If $n\geq5$, then $\dim W_2=\binom n2-1$ and
$\mathbf j\notin W_2$.
\end{lemma}
\begin{proof}
Suppose scalars $a_{ij}=a_{ji}$, indexed by distinct elements, satisfy
\[
 a_{ij}+a_{ik}+a_{jk}=c
 \qquad\text{for every triple }\{i,j,k\}.
\]
Fix distinct $k,l$ and put $b_i=a_{ik}-a_{il}$ for $i\notin\{k,l\}$.
Subtracting the equations for $\{i,j,k\}$ and $\{i,j,l\}$ gives
$b_i+b_j=0$. There are at least three available indices. For any three,
these equations give $b_i=-b_j=b_h=-b_i$, hence $2b_i=0$ and $b_i=0$.
Thus $a_{ik}=a_{il}$ for every allowable $i,k,l$. All coefficients on
edges incident with any fixed vertex are equal, and symmetry implies
that all $a_{ij}$ have one common value $a$. Consequently $c=3a=0$.

Applied with $c=0$, this identifies the kernel of the map
$(a_e)_e\mapsto\sum_e a_ev_e$ as the one-dimensional space of constant
coefficient vectors. Applied with $c=1$, it shows that $\mathbf j$ has
no preimage. Both assertions follow.
\end{proof}

\begin{theorem}\label{ULCD.TernaryKneser}
Let $n\geq7$ and $n\equiv5\pmod9$. Let $A$ be the adjacency matrix of
the Kneser graph on the triples of $\Omega$, so two triples are adjacent
if and only if they are disjoint. Then
\[
 \im(I+A)=W_2\oplus\langle\mathbf j\rangle,
 \qquad \im A=\bigl(W_2\oplus\langle\mathbf j\rangle\bigr)^\perp.
\]
Both spaces are LCD, and their dimensions are respectively
$\binom n2$ and $\binom n3-\binom n2$.
\end{theorem}
\begin{proof}
Let $M_s$ be the incidence matrix whose rows are the $s$-subsets and
whose columns are the triples, for $s=1,2$. For two triples whose
intersection has size $t$, inclusion and exclusion gives
\[
 \mathbf1_{t=0}=1-t+\binom t2-\binom t3.
\]
Since $\binom t3$ is the equality indicator for the two triples,
\begin{equation}\label{ULCD.TernaryIdentity}
 R:=I+A=J-M_1^{\mathsf T}M_1+M_2^{\mathsf T}M_2.
\end{equation}
For each point $i$, the sum of the vectors $v_e$ over the pairs $e$
containing $i$ is twice the point-incidence vector. Since two is
invertible in $\mathbb F_3$, the row space of $M_1$ is contained in
$W_2$. Equation~\eqref{ULCD.TernaryIdentity} therefore gives
$\im R\subseteq W_2+\langle\mathbf j\rangle$.

For a pair $e$ and a triple $T$, direct counting gives
\[
 (Rv_e)_T=\mathbf1_{e\subseteq T}
          +(n-5)\mathbf1_{e\cap T=\varnothing}.
\]
The second coefficient vanishes in $\mathbb F_3$. Thus $Rv_e=v_e$.
Also $R\mathbf j=(1+\binom{n-3}{3})\mathbf j=\mathbf j$:
writing $n=9m+5$, the binomial coefficient is
$3m(9m+2)(9m+1)/2$, which is divisible by three.
It follows that $R$ is the identity on $W_2+\langle\mathbf j\rangle$,
and hence that its image is exactly that space and $R^2=R$.

The matrix $R$ is symmetric. Lemma~\ref{ULCD.Projector} gives $\ker R=(\im R)^\perp$ and the orthogonal decomposition; both summands are LCD. Since $A=R-I$, its image is $\ker R$.
Lemma~\ref{ULCD.TernaryIncidence} now gives the two dimensions.
\end{proof}

\section{Intersection counts for subset and Hamming graphs}

Write $C_M=\im M$ for a symmetric matrix $M$ over $\mathbb F_p$. Let $A$
be the adjacency matrix of a regular simple graph on $N$ vertices, of degree
$\nu$, and put $B=J-I-A$. Following F017, call $C_A$ RLCD when
$C_A=C_{A+I}^{\perp}$, and CLCD when $C_A=C_B^{\perp}$.
Suppose ordered pairs of vertices are partitioned into classes, indexed
by $j$, on which both the adjacency indicator $a_j$ and the number $i_j$
of common neighbours are constant. Include the diagonal class among them.

\begin{theorem}\label{ULCD.Scalar}
Put $P_j=i_j+a_j$ and $\bar\nu=N-1-\nu$. Then:
\begin{enumerate}
\item $C_A$ is RLCD if and only if $P_j=0$ in $\mathbb F_p$ for every $j$.
\item $C_A\perp C_B$ if and only if $\nu-P_j=0$ for every $j$.
In this event both codes are LCD. They are complementary if and only if
$(\nu,\bar\nu)\ne(0,0)$ in the field; otherwise
$\mathbb F_p^N=C_A\perp\!\oplus C_B\perp\!\oplus\langle\mathbf1\rangle$.
\item $C_B$ is RLCD if and only if $\bar\nu-\nu+P_j=0$ for every $j$.
\end{enumerate}
\end{theorem}
\begin{proof}
The entries of $A^2+A$ are $P_j$. Thus their vanishing is exactly
$A(A+I)=0$. The matrices $-A$ and $I+A$ are then complementary symmetric
idempotents. Conversely the asserted duality implies that all their row
inner products vanish. Next
$AB=AJ-A-A^2=\nu J-A-A^2$, proving the orthogonality test.
The remaining claims in part~2 follow from Theorem~\ref{ULCD.Complement}, applied
to this actual pair of matrices; no condition on their complex spectra is
needed. Finally $B(B+I)=B(J-A)=\bar\nu J-AB$, which gives part~3 by the
same reflexive argument.
\end{proof}

\paragraph{Uniform subset graphs.}
Let $n\geq2k+1$, $k\geq2$ and $0\leq r<k$. The vertices of $\Gamma(n,k,r)$
are the $k$-subsets of an $n$-element set; adjacency means intersection of
size $r$. For two vertices meeting in $j$ elements, the common-neighbour
count is
\begin{equation}\label{ULCD.SubsetCounts}
 i_j=\sum_a\binom ja\binom{k-j}{r-a}^{\!2}
                  \binom{n-2k+j}{k-2r+a},\qquad 0\leq j\leq k.
\end{equation}
Here binomial coefficients with invalid lower argument are zero.
To prove the formula, partition the ground set into the intersection,
the two differences and the outside. A common neighbour chooses $a$
elements from the intersection, $r-a$ from each difference and the remaining
$k-2r+a$ outside. These choices are disjoint and exhaustive.
In particular $i_k=\nu=\binom kr\binom{n-k}{k-r}$ and
$P_j=i_j+\mathbf1_{j=r}$, so Theorem~\ref{ULCD.Scalar} applies directly.

For the Johnson graph $r=k-1$, the only nonzero counts off the diagonal
are $i_{k-2}=4$ and $i_{k-1}=n-2$. Thus RLCD holds exactly for $p=2$ and
odd $n$. Under those conditions $\nu=0$, so CLCD holds exactly when
$\binom nk$ is even. For $k\geq3$, complement orthogonality itself forces
$\nu=0$ from $P_0=0$;
then $P_{k-2}=4$ and $P_{k-1}=n-1$ force the same characteristic and parity.
For $k=2$, subtraction of the two off-diagonal tests yields $n=5$ in the
field, while $\nu=4$ gives $2(n-2)=4$, hence characteristic two and odd $n$.
The binary dimensions and minimum words follow from Theorems~\ref{ULCD.Johnson} and~\ref{ULCD.SimplexMin}, not from this count alone.

The domain qualification matters. F017 Corollary~1(1) does not state it
explicitly. For example $J(4,3)=K_4$ is RLCD in characteristic three.
More generally the adjacency matrix of $J(n,n-1)=K_n$, for $n\geq3$, is
$J-I$, so $(J-I)J=(n-1)J$: it is RLCD exactly when $p\mid(n-1)$.
Its complement is empty, so it is CLCD exactly when $J-I$ is invertible,
equivalently $p\nmid(n-1)$. Thus the unrestricted source exclusion is false.
Formula~\eqref{ULCD.SubsetCounts} itself remains valid for $n\geq k$ when only
realized intersections $\max(0,2k-n)\leq j\leq k$ are tested; the table
below retains its stated $n\geq2k+1$ domain.

For a Kneser graph ($r=0$), formula~\eqref{ULCD.SubsetCounts} reduces to
$i_j=\binom{n-2k+j}{k}$. If $n=2k+1$, then $i_0=0$ and $P_0=1$, so the
code is never RLCD. For $k\geq3$ it is not CLCD either: complement
orthogonality would force both $\nu=P_0=1$ and $\nu=i_1=0$.
For $k=2$ there is an exception. The Petersen graph has $(N,\nu)=(10,3)$,
$P_0=1$, $P_1=1$ and $P_2=3$. Complement orthogonality holds exactly
in characteristic two, where $\nu=1$ and $\bar\nu=0$; hence this code
is CLCD. F017 Corollary~1(2) states the exclusion for every $k$.
Its CLCD assertion requires this endpoint correction. Its RLCD exclusion
is valid for every $k\geq2$.

\begin{table}[ht]
\centering\small
\begin{tabular}{ccp{0.32\textwidth}p{0.32\textwidth}}
\toprule
$k$&$r$&RLCD&CLCD\\\midrule
2&0&$p=2$, $n\equiv3\pmod4$&$p=2$, $n\equiv1\pmod4$\\
3&0&$p=2$, $n\equiv1\pmod4$; $p=3$, $n\equiv5\pmod9$;
$p\geq5$, $n\equiv5\pmod p$&$p=2$, $n\equiv1\pmod4$;
$p=3$, $n\equiv2,8\pmod9$; $p\geq5$, $n\equiv5\pmod p$\\
4&0&$p=2$, $n\equiv7\pmod8$&$p=2$, $n\equiv3\pmod8$\\
3&1&$p=2$, $n\equiv0\pmod4$&$p=2$, $n\equiv0,1\pmod4$\\
4&1&$p=2$, $n\equiv1\pmod4$; $p=3$, $n\equiv6\pmod9$
&$p=2$, $n\equiv1\pmod8$; $p=3$, $n\equiv6\pmod9$\\
\bottomrule
\end{tabular}
\caption{Complete parameter conditions for the five families specified in
F017, with $n\geq2k+1$. Each entry is an if and only if statement.}
\label{ULCD.SubsetFamilies}
\end{table}

\paragraph{Completeness of the arithmetic classification.}
Here is a finite exact certificate for Table~\ref{ULCD.SubsetFamilies}, including
the exclusion of all other primes. For each fixed $(k,r)$, regard the
$P_j$ and $\nu-P_j$ as polynomials in $n$ over $\mathbb Q$ using
\eqref{ULCD.SubsetCounts}. For each list $f_0,\ldots,f_k$, the accompanying program
records polynomial coefficients $b_j$ and a monic polynomial $g$ such that
\[
 \sum_j b_jf_j=g,\qquad g\mid f_j\quad\text{for all }j.
\]
Both assertions are verified by exact polynomial arithmetic. Let $D$ be
the product of a common denominator of all coefficients of all $b_j$ and
a common denominator of all coefficients of all $f_j$.
For $p\nmid D$, the simultaneous congruences are equivalent to $g(n)=0$.
In the two $k=3,r=0$ lists, $g=n-5$; in all other lists, $g=1$.
The resulting exceptional primes, including primes subsequently excluded
by the finite tests, are
\[
\begin{array}{c|cc}
(k,r)&P_j&\nu-P_j\\\hline
(2,0)&2&2\\
(3,0)&2,3&2,3\\
(4,0)&2,3&2,3,17\\
(3,1)&2,3,5&2,3,23\\
(4,1)&2,3,5&2,3,5.
\end{array}
\]
For an exceptional prime choose the least power $p^a>k$. For $0\leq b\leq k$,
the identity $(1+X)^{p^a}=1+X^{p^a}$ over $\mathbb F_p$ proves
$\binom{m+p^a}{b}=\binom mb$ modulo $p$ for nonnegative $m$.
Consequently all the counts, the degree and $\binom nk$ have period $p^a$
modulo $p$ in the domain $n\geq2k+1$. Testing one admissible lift of every
residue modulo $p^a$ is therefore exhaustive, not a bounded experiment.
The program records these residues and applies the constant-line exclusion
in Theorem~\ref{ULCD.Scalar} to obtain CLCD. For $k=3,r=0,p\geq5,n\equiv5$,
the degree is zero but $\binom n3-1\equiv9\ne0$, so no further exception
occurs. This proves the table from its exact arithmetic certificate.

\paragraph{Hamming graphs.}
Let $H_k(d,q)$ have vertex set $\{0,\ldots,q-1\}^d$, with adjacency at
Hamming distance $k$, where $1\leq k\leq d$ and $q\geq2$. If two vertices
have distance $j$, their common-neighbour count is
\begin{equation}\label{ULCD.HammingCounts}
 i_j=\sum_a\binom{d-j}{a}(q-1)^a
 \frac{j!}{b!^2e!}(q-2)^e,
 \qquad b=j+a-k,\quad e=2k-j-2a,
\end{equation}
where terms with negative factorial arguments are omitted. In their $d-j$
equal coordinates, choose the $a$ positions where a common neighbour differs.
Among the $j$ different coordinates, let $b$ agree with the first vertex,
$b$ with the second, and $e$ with neither. The distance equations give the
stated values of $b,e$, with $2b+e=j$. Multinomial counting proves
\eqref{ULCD.HammingCounts}. Theorem~\ref{ULCD.Scalar} now applies with
$N=q^d$, $\nu=i_0=\binom dk(q-1)^k$ and $P_j=i_j+\mathbf1_{j=k}$.

\begin{corollary}\label{ULCD.HammingOddAlphabet}
For odd $q$, every binary distance-$k$ Hamming code is RLCD.
For the distance-two binary Hamming graph over $\mathbb F_3$, RLCD holds
whenever $d\equiv0\pmod3$.
\end{corollary}
\begin{proof}
Let $X_i$ act as the complete-graph adjacency matrix $J_q-I_q$ in coordinate
$i$ and as the identity in all others. Over $\mathbb F_2$, odd $q$ gives
$X_i^2=X_i$, and these matrices commute. The distance-$k$ adjacency matrix
is $A_k=\sum_{|S|=k}\prod_{i\in S}X_i$. Squaring cancels mixed terms in
characteristic two, so $A_k^2=A_k$ and $A_k(A_k+I)=0$.
For $q=2,k=2$, the only possible common-neighbour distances are $0,2,4$,
with counts $\binom d2$, $2(d-2)$ and $6$, respectively. If $3\mid d$,
these are $0,-1,0$ modulo three. This is precisely the RLCD test.
\end{proof}

For completeness, the Hamming spectrum used to compute the dimensions
comes from the same tensor construction. Over $\mathbb Q$, $J_q-I_q$
has eigenvalue $q-1$ on the constant vectors and $-1$ on their complement.
Choosing $t$ nonconstant tensor factors therefore gives
\[
 \lambda_t=\sum_i(-1)^i(q-1)^{k-i}\binom ti\binom{d-t}{k-i},
 \qquad m_t=\binom dt(q-1)^t.
\]
Lemma~\ref{ULCD.Rank} sums precisely those multiplicities for which
$p\nmid\lambda_t$. In particular the ternary $H_2(6,2)$ and $H_2(9,2)$
codes have dimensions $42$ and $342$. This rank step is justified by the
proved LCD property, rather than an assumption of diagonalizability after
reduction modulo $p$.

\begin{theorem}\label{ULCD.Cube}
Over $\mathbb F_2$, the distance-$k$ binary Hamming adjacency matrix satisfies
$A_k^2=\binom dk I$. For positive even $d$, the code of the ordinary cube
$Q_d=H_1(d,2)$ is self-dual of dimension $2^{d-1}$. Its complement adjacency
matrix is invertible, so the code is not the dual of the complement code.
\end{theorem}
\begin{proof}
Index coordinates by $\mathbb F_2^d$. Write $T_s$ for translation by $s$.
Then $A_k=\sum_{\wt(s)=k}T_s$. The translations commute, each satisfies
$T_s^2=I$, and cross terms cancel in the square, proving the first assertion.
Write $M_t$ for the adjacency matrix of the ordinary cube $Q_t$.
Order cube vertices by their last coordinate. If $d$ is even,
\[
 M_d=\begin{pmatrix}M_{d-1}&I\\I&M_{d-1}\end{pmatrix},\qquad
 M_{d-1}^2=I.
\]
The kernel consists exactly of $(x,M_{d-1}x)$, so both nullity and rank are
$2^{d-1}$. Since $M_d^2=0$ and $M_d$ is symmetric, its image equals its
orthogonal complement. On the other hand, with $B=J+I+M_d$ we have
$J^2=2^dJ=0$ and $JM_d=M_dJ=dJ=0$, so $B^2=I$.
\end{proof}

The conclusion of F017 displays the complement-dual identity for
$Q_6,Q_8,Q_{10}$. Direct reading of the formula confirms the complement bar;
Theorem~\ref{ULCD.Cube} disproves that identity and supplies the corrected
self-duality statement. The two remaining Hamming examples in that
conclusion are instances of Corollary~\ref{ULCD.HammingOddAlphabet} and are retained.

\section{Finite certificates for modules and automorphisms}

The following criteria explain what the finite data in the applications
must certify. They use ordinary linear algebra and group actions. A list
of computed ranks is evidence only after its matrices have been identified
with the specified code, action or quotient.

\begin{lemma}\label{ULCD.SquareZeroIncidence}
Let $R$ be a matrix over $F$ such that $RR^{\mathsf T}=0$. Then
$Q=R^{\mathsf T}R$ is symmetric, $Q^2=0$, and $\im Q$ is self-orthogonal.
If $Rj=0$ and $j^{\mathsf T}j=0$, then $\im Q+Fj$ is self-orthogonal too.
\end{lemma}
\begin{proof}
We have $Q^2=R^{\mathsf T}(RR^{\mathsf T})R=0$, and hence
$\langle Qx,Qy\rangle=x^{\mathsf T}Q^2y=0$.
Also $Qj=R^{\mathsf T}Rj=0$, so $j$ is orthogonal to $\im Q$.
The remaining pairing is the assumed norm of $j$.
\end{proof}

\begin{theorem}\label{ULCD.MatrixAlgebra}
Let $d\geq1$, and let $V=F^d$ carry a representation of a group generated by
$g_1,\ldots,g_s$, with matrices $M_1,\ldots,M_s$.
If $d^2$ products of these matrices, including the empty product if needed,
are linearly independent in $\operatorname{End}_F(V)$, then $V$ is
absolutely irreducible.
It suffices instead to exhibit a rank-one matrix $P=vu$ in the generated
unital algebra such that the columns $Mv$ and rows $uN$, as $M,N$ run
over products of the generators, each span a space of dimension $d$.
\end{theorem}
\begin{proof}
The first hypothesis says that the generated algebra is
$\operatorname{End}_F(V)$. Every group-invariant subspace is stable under
this algebra. If it contains a nonzero vector, an endomorphism can send
that vector to any specified vector of $V$, so the subspace is $V$.
Linear independence of the $d^2$ matrices persists under every field
extension: a nonzero full-size minor stays nonzero. The same argument
therefore proves irreducibility after extension.

For the second test, select products $M_i,N_j$ so that the $d$ columns
$M_iv$ and the $d$ rows $uN_j$ are bases. Their $d^2$ outer products
form a basis of the matrix space: changing either chosen basis to the
standard basis is invertible. Each outer product is $M_iPN_j$ and belongs
to the generated algebra. The first argument applies.
\end{proof}

\begin{theorem}\label{ULCD.Socle}
Let $V$ be a finite-dimensional $FG$-module, with an invariant filtration
$0=V_0<V_1<\cdots<V_s=V$ whose factors $D_i=V_i/V_{i-1}$ are simple.
Every simple submodule of $V$ is the image of an injective homomorphism
from one of the $D_i$. Consequently, if a semisimple submodule $N$ contains
the image of every homomorphism $D_i\to V$, then $N=\operatorname{Soc}(V)$.
For this containment it suffices to check every map in a vector-space
basis of each $\operatorname{Hom}_G(D_i,V)$.
\end{theorem}
\begin{proof}
Let $S\subseteq V$ be simple, and choose the least $i$ with $S\subseteq V_i$.
The invariant subspace $S\cap V_{i-1}$ is either zero or $S$; minimality
excludes the latter. Projection therefore embeds $S$ in the simple module
$D_i$, and is an isomorphism. Its inverse followed by inclusion in $V$
is the required homomorphism. Under the containment hypothesis every
simple submodule lies in $N$, so their sum lies in $N$. Conversely a
semisimple submodule is a sum of simple submodules and lies in that sum.
Finally an arbitrary homomorphism is a linear combination of basis maps;
its image lies in the sum of their images.
\end{proof}

In row coordinates the homomorphisms are exactly the solutions of
$XP_g=D_gX$ for each chosen generator. Solving the whole homogeneous
system provides the required basis. Checking a selected map would not
give the conclusion. Invariant prefixes are checked directly on each
generator; simplicity of their factors can be certified by
Theorem~\ref{ULCD.MatrixAlgebra}. Thus the criterion does not require a
classification of all irreducible representations of $G$.

\begin{theorem}\label{ULCD.Refinement}
Let $\mathcal S$ be a collection of distinct subsets of a finite set
$\Omega$. Form its bipartite incidence graph, with the two parts initially
coloured differently. Refine a partition by recording, for every vertex,
its old cell and its numbers of neighbours in every old cell, until no
cell splits. Choose $x_1,\ldots,x_t\in\Omega$ successively, separating
$x_i$ from its current cell and refining again. If that cell has size $b_i$
before separation, and all vertices of $\Omega$ have different colours
at the end, then
\[
 |\operatorname{Aut}(\Omega,\mathcal S)|\leq\prod_{i=1}^t b_i.
\]
In particular this bounds the coordinate-permutation group of any code
whose complete set of minimum supports is $\mathcal S$.
\end{theorem}
\begin{proof}
An incidence automorphism preserves the initial parts. Whenever it
preserves every old cell it also preserves the neighbour counts defining
the refined cells. After separating $x_1,\ldots,x_{i-1}$, the stabilizer
of those points therefore preserves every current cell. Its orbit of
$x_i$ has size at most $b_i$. At the last stage this stabilizer fixes
every point of $\Omega$ and hence every subset in $\mathcal S$.
It is trivial. Repeated orbit--stabilizer gives the inequality.
A coordinate permutation preserving a code preserves word weights and
therefore its complete minimum-support collection, proving the last claim.
\end{proof}

To obtain an equality, the application supplies actual permutations
preserving the code and verifies that they generate a subgroup of order
$\prod b_i$. The upper bound concerns the full group, independently of
how those permutations were discovered. For a graph with adjacency matrix
$A$, every graph automorphism preserves $\im A$; thus the additional
inclusions $G\leq\operatorname{Aut}(\Gamma)\leq\operatorname{Aut}(C_A)$
identify the full graph group whenever both outer groups have that order.

\section{Weights and dual weights}

The averaging theorem below limits a distance search when the code has
a transitive permutation group. A complete distribution requires either
enumeration of the whole code or another proof. In particular a count of
messages of bounded weight is not a whole-code weight distribution.

\begin{theorem}[MacWilliams identity]\label{ULCD.MacWilliams}
Let $C\leq\mathbb F_q^n$ have dimension $k$, and let $A_i$ be its number
of words of weight $i$. Write $A_j^\perp$ for the corresponding numbers
in the orthogonal dual. Then
\[
 A_j^\perp=q^{-k}\sum_{i=0}^n A_i
 \sum_{s=0}^j(-1)^s(q-1)^{j-s}
                 \binom is\binom{n-i}{j-s}.
\]
\end{theorem}
\begin{proof}
Choose a nontrivial additive character
$\chi:\mathbb F_q\to\mathbb C^\times$. Such a character is obtained
from a nonzero $\mathbb F_p$-linear functional to $\mathbb F_p$ and a
primitive $p$th root of unity, where $p$ is the characteristic.
For a nontrivial character on a finite additive group its sum is zero:
translation by an element on which it is not one multiplies the sum by
that value and also leaves it unchanged.
For fixed $x\in\mathbb F_q^n$, the map $c\mapsto c\cdot x$ is either
zero on $C$ or is onto $\mathbb F_q$, with equally sized fibres.
It follows that
\[
 q^{-k}\sum_{c\in C}\chi(c\cdot x)=\mathbf1_{x\in C^\perp}.
\]
Multiply by $z^{\mathrm{wt}(x)}$ and sum over all $x$. The finite sums
may be interchanged. For a fixed $c$ of weight $i$, the coordinate factors
in the inner sum are $1+(q-1)z$ at its zero entries and $1-z$ at its
nonzero entries. Hence
\[
 \sum_j A_j^\perp z^j
 =q^{-k}\sum_i A_i(1-z)^i(1+(q-1)z)^{n-i}.
\]
Taking the coefficient of $z^j$ proves the formula.
\end{proof}

This is a classical identity, included here to make its repeated use
explicit. For a matrix with $k$ independent rows, the map from its
$q^k$ coefficient vectors to its row space is bijective. Enumerating
each coefficient vector once therefore gives the complete $A_i$.
The displayed identity then determines the dual distribution. Integrality,
nonnegativity and totals provide consistency checks; these checks alone
would not prove that a proposed input distribution belongs to the code.
\section{Transitive distance searches and rank}

Let $C\leq \mathbb F_p^n$ be a nonzero linear code of dimension $k$.
An information set is a set $I$ of $k$ coordinates such that the restriction
$C\longrightarrow\mathbb F_p^I$ is an isomorphism. Gaussian elimination
produces a generator matrix $G$ whose columns indexed by $I$ form the identity.
For this choice of $G$, the message of $c=uG$ is exactly $c|_I=u$.

\begin{theorem}\label{ULCD.TransitiveDistance}
Suppose that a group $T$ of coordinate permutations preserves $C$ and is
transitive on the $n$ coordinates. Fix an information set $I$. For every
$c\in C$ of weight $w$, some $t\in T$ satisfies
\[
 \wt((tc)|_I)\leq\left\lfloor\frac{kw}{n}\right\rfloor.
\]
Consequently, for any positive integer $D$, the following finite data prove
that $C$ has minimum distance $D$:
\begin{enumerate}
\item a codeword of weight $D$;
\item verification that each nonzero message of weight at most
$\lfloor k(D-1)/n\rfloor$ gives a codeword of weight at least $D$.
\end{enumerate}
\end{theorem}
\begin{proof}
The permutation group $T$ is finite. For each fixed coordinate $x$ and each
$y$, transitivity implies that exactly $|T|/n$ elements send $x$ to $y$:
each nonempty fibre is a coset of a point stabilizer. Count pairs $(t,x)$
with $x\in\operatorname{supp}(c)$ and $tx\in I$. This gives
\[
 \sum_{t\in T}\wt((tc)|_I)=\frac{|T|kw}{n}.
\]
At least one summand is no larger than the average, proving the assertion.
If a nonzero codeword had weight $w<D$, its image under the corresponding
permutation would have nonzero message of weight at most
$\lfloor kw/n\rfloor\leq\lfloor k(D-1)/n\rfloor$. Its codeword weight is
still $w$, contradicting item~2. Item~1 proves the reverse bound.
\end{proof}

\begin{corollary}\label{ULCD.MinimumEnumeration}
Under the same hypotheses, suppose the minimum distance is $D$ and $S$ is
a $T$-invariant set of words of weight $D$ in $C$. If every word of weight
$D$ obtained from a message of weight at most $\lfloor kD/n\rfloor$
belongs to $S$, then $S$ is the complete set of minimum words.
\end{corollary}
\begin{proof}
Apply Theorem~\ref{ULCD.TransitiveDistance} to an arbitrary minimum word. Its resulting
representative is in $S$, and $T$-invariance returns the original word to $S$.
The opposite inclusion is part of the hypothesis.
\end{proof}

The number of messages in the distance test, including zero, is
\[
 \sum_{i=0}^{\lfloor k(D-1)/n\rfloor}\binom{k}{i}(p-1)^i.
\]
The theorem concerns this complete finite set. It does not treat a random
sample as a certificate, and its restricted weight counts are not the weight
distribution of the whole code. Transitivity is essential. For example, the
binary code generated by $(1,0,0,0)$ has $n=4$, $k=1$, $d=1$ and information
set $\{1\}$. A putative lower-bound test for $D=2$ enumerates only the zero
message and would give a false conclusion if the transitivity hypothesis
were removed.

For a proposed invariant set of minimum words, the hypothesis in
Corollary~\ref{ULCD.MinimumEnumeration} can also be checked by counting.
If its intersection with the restricted message set has the same size
as the complete restricted minimum set, inclusion and equality of these
finite sizes prove equality of sets. No disjointness or orbit count is
assumed without checking it.

\begin{lemma}\label{ULCD.Rank}
Let $A$ be a symmetric integer matrix with integral eigenvalues
$\lambda_1,\ldots,\lambda_s$ of multiplicities $m_1,\ldots,m_s$ over
$\mathbb Q$. Suppose that the code $\im\bar A$ over $\mathbb F_p$ is LCD,
where a bar denotes reduction modulo $p$. Then
\[
 \rank\bar A=\sum_{p\nmid\lambda_i}m_i.
\]
\end{lemma}
\begin{proof}
Symmetry gives $(\im\bar A)^\perp=\ker\bar A$. The LCD hypothesis and
rank--nullity therefore give the invariant direct sum
$\mathbb F_p^n=\ker\bar A\oplus\im\bar A$. The restriction of $\bar A$
to its image is injective, hence invertible. Its characteristic polynomial
has nonzero constant coefficient. Thus the multiplicity of zero in the
characteristic polynomial of $\bar A$ is exactly $\dim\ker\bar A$.
Reduction of the integer characteristic polynomial gives
$\prod_i(X-\bar\lambda_i)^{m_i}$, whose zero multiplicity is
$\sum_{p\mid\lambda_i}m_i$. Subtracting from $n$ proves the formula.
\end{proof}


\paragraph{Required finite arithmetic.}
The polynomial identities and complete residue lists for
Table~\ref{ULCD.SubsetFamilies} are certified by
\path{check_association_counts.py}. The proof explains why these finite
identities settle every prime and every admissible parameter, rather than
inferring an infinite statement from sampled graphs. The applications
specify the further code and group certificates they require. No finite
distance, irreducibility or full automorphism assertion follows merely
from an adjacency spectrum or a congruence test.
\end{document}
